\chapter{Banda de partícula, anti-sección y virtualidad}
\label{chap:banda-particula-virtual}
\index{partícula!banda de ruta}
\index{antipartícula!anti-sección orientada}
\index{partícula virtual}

La persistencia de una ruta produce una partícula antes de que se evalúe su
masa. Esta prioridad permite distinguir tres objetos:

\begin{enumerate}
\item la \emph{sección orientada}, que conserva la dirección de la ruta;
\item la \emph{banda}, que reúne una sección y su conjugada;
\item la \emph{persistencia}, que decide si la firma alcanza o no una rama
global.
\end{enumerate}

\begin{definition}[Palabras dodecafásicas y conjugación material]
\label{def:palabras-cm}
Sea
\[
 \mathscr T_{12}^{\rm word}:=\{-1,0,1\}^{12}.
\]
Para \(\mathbf t=(t_1,\ldots,t_{12})\) escribimos
\[
 \operatorname{rev}(\mathbf t)
 =(t_{12},\ldots,t_1),
 \qquad
 C_M(\mathbf t):=-\operatorname{rev}(\mathbf t).
\]
La inversión \(C_M\) cambia simultáneamente el sentido de lectura y la
orientación trítica.
\end{definition}

\begin{theorem}[Involución material y sector fijo]
\label{thm:involucion-cm-censo}
La aplicación \(C_M\) es una involución de
\(\mathscr T_{12}^{\rm word}\) y
\[
 \left|\operatorname{Fix}(C_M)\right|=3^6=729.
\]
\end{theorem}

\begin{proof}
La reversión y el cambio de signo conmutan y ambos tienen cuadrado igual a
la identidad; por tanto \(C_M^2=I\). Una palabra es fija si y sólo si
\[
 t_i=-t_{13-i},\qquad 1\leq i\leq6.
\]
Las seis primeras coordenadas son libres y las seis últimas quedan
determinadas por ellas. Cada coordenada libre admite tres valores, de modo
que el sector fijo tiene cardinal \(3^6\).
\end{proof}

Esta conjugación finita no se obtiene cambiando únicamente el signo visible:
invierte también la orientación del registro.

\section{La banda de ruta}

Sea \(\mathscr T_{12}^{\rm word}\) el espacio de palabras de la
\cref{def:palabras-cm}. Añadimos una orientación de sección
\(\chi\in\{-1,+1\}\) y definimos
\[
(\mathbf t,\chi)\sim(C_M\mathbf t,-\chi).
\]

\begin{definition}[Banda dodecafásica de ruta]
\label{def:banda-dodecafasica-ruta}
La banda de la sección orientada \((\mathbf t,\chi)\) es la clase
\[
[(\mathbf t,\chi)]_M
=\{(\mathbf t,\chi),(C_M\mathbf t,-\chi)\},
\]
y el espacio de bandas es
\[
\mathcal B_{12}
:=
\bigl(\mathscr T_{12}^{\rm word}\times\{-1,+1\}\bigr)/{\sim}.
\]
\end{definition}

\begin{proposition}[Libertad de la conjugación orientada]
\label{prop:accion-cm-orientada-libre}
La acción
\[
(\mathbf t,\chi)\longmapsto(C_M\mathbf t,-\chi)
\]
es libre. En consecuencia,
\[
|\mathcal B_{12}|=3^{12}.
\]
\end{proposition}

\begin{proof}
Un punto fijo exigiría \(\chi=-\chi\), imposible para
\(\chi\in\{-1,+1\}\). Todas las órbitas tienen dos elementos y el dominio
posee \(2\cdot3^{12}\).
\end{proof}

La banda no borra la orientación: la conserva como doble cubierta. La
partícula y la antipartícula son las dos secciones orientadas de la misma
banda.

\section{Masa par y carga impar}

Sea \(E_0:\mathscr T_{12}^{\rm word}\to\mathbb R\) un carácter aditivo de
ruta. Sus componentes par e impar respecto de \(C_M\) son
\[
E_+(\mathbf t)
=\frac12\bigl(E_0(\mathbf t)+E_0(C_M\mathbf t)\bigr),
\]
\[
E_-(\mathbf t)
=\frac12\bigl(E_0(\mathbf t)-E_0(C_M\mathbf t)\bigr).
\]
Entonces
\[
E_+(C_M\mathbf t)=E_+(\mathbf t),
\qquad
E_-(C_M\mathbf t)=-E_-(\mathbf t).
\]

\begin{definition}[Carácter de banda]
\label{def:caracter-banda}
Para una sección orientada definimos
\[
E_{\rm band}(\mathbf t,\chi)
=E_+(\mathbf t)+\chi E_-(\mathbf t),
\qquad
m(\mathbf t,\chi)
=m_\circ\exp\bigl(E_{\rm band}(\mathbf t,\chi)\bigr),
\]
donde \(m_\circ\) es la unidad interna de comparación de la familia.
\end{definition}

\begin{theorem}[Principio de conjugación de masa]
\label{thm:principio-conjugacion-masa-banda}
El carácter y la masa descienden al cociente de bandas:
\[
E_{\rm band}(C_M\mathbf t,-\chi)
=E_{\rm band}(\mathbf t,\chi),
\]
\[
\boxed{
m(C_M\mathbf t,-\chi)=m(\mathbf t,\chi).}
\]
Los invariantes pares pertenecen a la banda; los invariantes impares
distinguen la orientación de sus dos secciones.
\end{theorem}

\begin{proof}
Por las paridades anteriores,
\[
\begin{aligned}
E_{\rm band}(C_M\mathbf t,-\chi)
&=E_+(C_M\mathbf t)-\chi E_-(C_M\mathbf t)\\
&=E_+(\mathbf t)+\chi E_-(\mathbf t).
\end{aligned}
\]
La igualdad de masas se obtiene al aplicar la exponencial.
\end{proof}

En particular, una forma lineal con pesos simétricos es impar, mientras una
forma cuadrática compatible con la reversión es par. La carga y la
orientación cambian; la masa, construida sobre la banda, permanece.

\begin{corollary}[Electrón y positrón]
\label{cor:electron-positron-banda}
En la ecuación electrónica de banda, el par
\((n,\chi)=(-22,+1)\) y su conjugado \((+22,-1)\) producen el mismo término
\(\chi n=-22\). Por tanto, electrón y positrón poseen la misma masa y
orientaciones de carga opuestas.
\end{corollary}

\section{Cierre finito del atlas bajo anti-sección}

En las ochenta y una rutas CT--108 aparecen siete formas tríticas. El censo
exhaustivo distingue tres operaciones:
\[
\mathbf t\longmapsto-\mathbf t,\qquad
\mathbf t\longmapsto\operatorname{rev}(\mathbf t),\qquad
\mathbf t\longmapsto-\operatorname{rev}(\mathbf t).
\]

\begin{proposition}[Cierre de las formas CT--108]
\label{prop:cierre-antiseccion-ct108}
El conjunto de formas tríticas de CT--108 no es estable bajo signo puro ni
bajo reversión pura, y es estable bajo \(C_M=-\operatorname{rev}\). Sus
órbitas poseen multiplicidades
\[
\boxed{81=54+9+9+9.}
\]
La órbita de multiplicidad \(54\) es auto-conjugada; las otras tres
intercambian pares de formas.
\end{proposition}

\begin{proof}
La forma fija es
\[
\texttt{+++---+++---}
\]
y posee multiplicidad \(54\). Las tres parejas
\[
\begin{array}{c}
\texttt{+++--0+++--0}\ /\ \texttt{0++---0++---},\\
\texttt{+++-0-+++-0-}\ /\ \texttt{+0+---+0+---},\\
\texttt{+++0--+++0--}\ /\ \texttt{++0---++0---}
\end{array}
\]
poseen multiplicidad total \(9\) cada una. La aplicación
\(-\operatorname{rev}\) fija la primera y permuta los miembros de cada
pareja. El cambio de signo y la reversión por separado producen palabras
fuera de esta lista. La suma de multiplicidades es
\(54+9+9+9=81\).
\end{proof}

El resultado no identifica el sector \(54\) con otro objeto de cardinal
homónimo. Demuestra qué conjugación respeta la taxonomía concreta de las
rutas y por qué partícula y antipartícula comparten banda.

\section{Persistencia finita y partícula virtual}

Sea
\[
\cdots\xrightarrow{\rho_{M+1,M}}S_M
\xrightarrow{\rho_{M,M-1}}S_{M-1}\xrightarrow{}\cdots
\]
el sistema de estados admisibles a profundidades finitas. Para
\(s_N\in S_N\), definimos el conjunto de prolongaciones
\[
\mathcal P_M(s_N)
:=
\{s_M\in S_M:\rho_{M,N}(s_M)=s_N\},
\qquad M\geq N.
\]
Su profundidad de vida es
\[
\mathcal L(s_N)
:=\sup\{M\geq N:\mathcal P_M(s_N)\neq\varnothing\}
\in\mathbb N\cup\{\infty\}.
\]

\begin{definition}[Persistencia y virtualidad]
\label{def:persistencia-virtualidad}
Un prefijo \(s_N\) es \emph{persistente} si pertenece a una sección global
\((s_M)_{M\geq N}\) compatible. Es \emph{virtual finito} si
\(\mathcal L(s_N)<\infty\), existe una prolongación hasta
\(\mathcal L(s_N)\) y no existe ninguna a la profundidad siguiente.
\end{definition}

\begin{theorem}[Criterio de persistencia]
\label{thm:criterio-persistencia-virtual}
Supóngase que cada estado posee un número finito de prolongaciones inmediatas
y que las restricciones son coherentes. Entonces:
\[
\mathcal L(s_N)=\infty
\quad\Longleftrightarrow\quad
s_N\text{ pertenece a una sección global}.
\]
Si \(\mathcal L(s_N)=L<\infty\), el prefijo transporta ruta, memoria,
registro y firma hasta \(L\), pero no define una partícula persistente.
\end{theorem}

\begin{proof}
Una sección global proporciona una prolongación a toda profundidad, de modo
que \(\mathcal L=\infty\). Recíprocamente, las prolongaciones de \(s_N\)
forman un árbol enraizado, infinito y finitamente ramificado. Por el lema de
Kőnig contiene una rama infinita, cuyos vértices son compatibles por
construcción y definen una sección global.

Si \(\mathcal L=L<\infty\), existe un estado a profundidad \(L\) y ninguno
a \(L+1\). Los lectores definidos sobre prefijos pueden evaluarlo hasta
\(L\); la ausencia de prolongación impide convertir esa firma finita en una
sección del límite inverso.
\end{proof}

La partícula virtual posee así una definición interna. No es una partícula
``menos real'' ni una violación breve de una ley de conservación. Es una
ruta cuya memoria influye durante una profundidad finita y cuya genealogía
se extingue antes de estabilizarse. El contraste con propagadores y estados
intermedios de la física convencional se realiza después de identificar qué
lector representa esas profundidades.

\section{Cadena rectora}

Los resultados de este capítulo fijan la dirección completa:
\[
\boxed{
\text{extensión}\to
\text{ruta}\to
\text{registro}\to
\text{firma}\to
\text{banda orientada}\to
\text{partícula}\to
\text{masa}.}
\]
La antipartícula es la segunda sección de la banda. La virtualidad es una
persistencia finita. Ninguna de estas nociones se define a partir de una
masa observada.
